\vspace{0.5em}\noindent\textbf{Internship Application
Platform}\par

\textbf{An AI-integrated internship application management platform
deployed on AWS}

\vspace{0.5em}\noindent\textbf{Project Summary}\par

Internship Application Platform is a centralized platform that helps
students, candidates, and companies manage the internship recruitment
process in one system.

For Candidates, the system supports profile creation, internship
opportunity discovery, application status tracking, CV and document
management, direct communication with HR, and AI-based evaluation of the
fit between a CV and a Job Description.

For HR teams and companies, the system supports company profile
management, job posting, applicant review, recruitment status updates,
direct communication with Candidates, and AI-assisted CV analysis,
comparison, and ranking.

The project is designed as a multi-service architecture, including:

\begin{itemize}
\tightlist
\item
  A React and Vite frontend.
\item
  A FastAPI REST backend.
\item
  A chat service using Node.js, Socket.IO, Redis, and DynamoDB.
\item
  An AI service for CV analysis, Job Description analysis, and candidate
  ranking.
\item
  PostgreSQL for business data.
\item
  Amazon S3 for CV and document storage.
\item
  Docker and Kubernetes for standardized deployment environments.
\item
  Observability for metrics, logs, and distributed traces.
\item
  GitHub Actions for CI/CD and security checks.
\end{itemize}

The production architecture is planned for AWS using Amazon EKS, Amazon
ECR, Amazon RDS for PostgreSQL, Amazon DynamoDB, Amazon ElastiCache,
Amazon S3, Amazon CloudFront, Application Load Balancer, and suitable
monitoring services.

\vspace{0.5em}\noindent\textbf{Problem Statement}\par

\vspace{0.5em}\noindent\textbf{Current Problems}\par

During the internship search process, students often manage application
information with spreadsheets, personal notes, email, or several
separate platforms. This approach creates several problems:

\begin{itemize}
\tightlist
\item
  Job and company information is scattered.
\item
  It is difficult to track the status of each application.
\item
  CVs, transcripts, certificates, and related documents are stored in
  multiple places.
\item
  Candidates have limited support for evaluating how well their CV
  matches a Job Description.
\item
  Communication with HR is not directly connected to the application
  record.
\item
  Candidates can easily miss deadlines, interview schedules, or document
  requests.
\end{itemize}

For HR teams and companies, manual recruitment workflows also create
limitations:

\begin{itemize}
\tightlist
\item
  CV reading and classification takes a significant amount of time.
\item
  It is difficult to compare multiple candidates using consistent
  criteria.
\item
  Recruitment status may not be updated consistently.
\item
  Candidate data may be stored insecurely.
\item
  Communication between HR and Candidates is scattered across multiple
  channels.
\item
  There is no centralized tool for tracking recruitment activity and
  effectiveness.
\end{itemize}

\vspace{0.5em}\noindent\textbf{Proposed Solution}\par

Internship Application Platform provides a centralized system for both
Candidates and HR.

Candidates can:

\begin{itemize}
\tightlist
\item
  Register and manage their accounts.
\item
  Update personal profiles.
\item
  Search and view job information.
\item
  Submit applications.
\item
  Track recruitment status.
\item
  Upload and manage CVs, certificates, and transcripts.
\item
  Communicate directly with HR.
\item
  Use AI to analyze CV fit against a job.
\end{itemize}

HR teams and companies can:

\begin{itemize}
\tightlist
\item
  Create and manage company profiles.
\item
  Post and edit job openings.
\item
  View submitted applications.
\item
  Update candidate recruitment status.
\item
  Communicate directly with Candidates.
\item
  Use AI to extract information, calculate fit scores, and support
  candidate ranking.
\end{itemize}

Long-running tasks such as document reading, CV analysis, and reranking
are processed asynchronously through workers so that the main API is not
blocked during processing.

\vspace{0.5em}\noindent\textbf{Benefits and Value}\par

For Candidates, the platform helps:

\begin{itemize}
\tightlist
\item
  Manage the full application process in one place.
\item
  Reduce the risk of losing documents and job information.
\item
  Track the status of each application clearly.
\item
  Identify relevant skills and missing skills.
\item
  Communicate with companies more conveniently.
\end{itemize}

For HR teams, the platform helps:

\begin{itemize}
\tightlist
\item
  Reduce manual CV processing time.
\item
  Centralize candidate information and recruitment status.
\item
  Standardize the resume evaluation process.
\item
  Support data-driven decision-making.
\item
  Improve recruitment tracking and coordination.
\end{itemize}

For the development team, the project provides an opportunity to apply
full-stack development, Cloud architecture, AI integration, realtime
communication, Kubernetes, CI/CD, observability, and system security
knowledge.

\vspace{0.5em}\noindent\textbf{Solution Architecture}\par

The original proposal targeted a multi-service AWS architecture with
Kubernetes for containerized services, managed databases, object
storage, realtime chat, observability, and CI/CD. In the final
implementation, I kept that direction and adjusted two important areas
after completing the deployment work:

\begin{itemize}
\tightlist
\item
  The React frontend is no longer a Kubernetes workload. It is built
  with Vite, stored in a private S3 bucket, and delivered by CloudFront.
\item
  Backend, chat, worker, outbox dispatcher, and the SageMaker adapter
  remain on EKS because they are long-running processes or need
  Kubernetes rollout, health check, and scaling controls.
\end{itemize}

\vspace{0.5em}\noindent\textbf{Final implemented
architecture}\par

\{\{\textless{} mermaid \textgreater\}\} graph LR User{[}``Candidate /
HR browser''{]} CF{[}``Amazon
CloudFrontdhm2rz5nmsibj.cloudfront.net''{]} S3Frontend{[}``Private S3
frontend bucketinternship-prod-frontend-account-redacted''{]}
ALB{[}``Application Load Balancerinternet-facing''{]} EKS{[}``Amazon
EKSinternship-prod / namespace internship''{]} Backend{[}``FastAPI
backendDeployment/backend :8000''{]} Chat{[}``Node.js Socket.IO
chatDeployment/chat-service :3000''{]} Dispatcher{[}``Outbox
dispatcherDeployment/backend-outbox-dispatcher''{]}
Worker{[}``Processing workerDeployment/backend-processing-worker''{]}
AI{[}``AI service adapterDeployment/ai-service :8010''{]} RDS{[}``Amazon
RDS PostgreSQLinternship-prod-postgres''{]} Redis{[}``Amazon ElastiCache
Redisinternship-prod-redis''{]} DDBChat{[}``DynamoDB chat
tablesChatUsers / ChatGroups / ChatMessages''{]} S3Uploads{[}``S3
uploads and archive bucketinternship-prod-uploads-account-redacted''{]}
SQS{[}``Amazon SQSinternship-prod-outbox''{]} DLQ{[}``SQS
DLQinternship-prod-outbox-dlq''{]} Lambda{[}``AWS
Lambdainternship-outbox-handler''{]} Dedupe{[}``DynamoDB dedupe
tableInternshipLambdaEventDedupe''{]} SES{[}``Amazon SES''{]}
SageMaker{[}``SageMaker endpointinternship-qwen3-4b''{]}
GitHub{[}``GitHub Actions OIDC''{]} ECR{[}``Amazon ECR images''{]}

\begin{verbatim}
User --> CF
CF -->|"Default *"| S3Frontend
CF -->|"/api/*"| ALB
CF -->|"/chat/*"| ALB
CF -->|"/socket.io/*"| ALB
ALB --> Backend
ALB --> Chat
Backend --> RDS
Backend --> S3Uploads
Backend --> Dispatcher
Backend --> Worker
Chat --> Redis
Chat --> DDBChat
Worker --> AI
AI --> SageMaker
Dispatcher --> SQS
SQS --> Lambda
SQS --> DLQ
Lambda --> Dedupe
Lambda --> S3Uploads
Lambda --> SES
GitHub --> ECR
GitHub --> EKS
GitHub --> S3Frontend
GitHub --> CF
\end{verbatim}

\{\{\textless{} /mermaid \textgreater\}\}

\vspace{0.5em}\noindent\textbf{Architecture explanation}\par

CloudFront is the single public entry point for users. Static assets use
the default CloudFront behavior and are read from the private frontend
S3 bucket through CloudFront Origin Access Control. Dynamic API and
realtime routes are sent to the public Application Load Balancer:

\begin{itemize}
\tightlist
\item
  \texttt{/\allowbreak{}api/\allowbreak{}*} is rewritten by the ALB Ingress and routed to the
  FastAPI backend.
\item
  \texttt{/\allowbreak{}chat/\allowbreak{}*} is rewritten and routed to the chat service.
\item
  \texttt{/\allowbreak{}socket.\allowbreak{}io/\allowbreak{}*} is routed to the chat service for Socket.IO
  transport.
\end{itemize}

The EKS cluster hosts only services that need long-running runtime
behavior: backend, chat, outbox dispatcher, processing worker, and the
AI adapter. PostgreSQL is the transactional database for users, jobs,
applications, workflow state, async processing jobs, outbox records, and
idempotency records. DynamoDB stores permanent chat entities and Lambda
event deduplication state. Redis is used only for Socket.IO pub/sub
between chat replicas. SQS decouples committed business events from
notification processing, and Lambda performs short event-driven work
such as deduplication, S3 archive, and SES email delivery.

\vspace{0.5em}\noindent\textbf{Component responsibility
table}\par

{
\begingroup
\small
\setlength{\tabcolsep}{3pt}
\renewcommand{\arraystretch}{1.15}
\begin{longtable}{@{}
  >{\raggedright\arraybackslash}p{(\linewidth - 4\tabcolsep) * \real{0.3333}}
  >{\raggedright\arraybackslash}p{(\linewidth - 4\tabcolsep) * \real{0.3333}}
  >{\raggedright\arraybackslash}p{(\linewidth - 4\tabcolsep) * \real{0.3333}}@{}}
\toprule\noalign{}
\begin{minipage}[b]{\linewidth}\raggedright
Component
\end{minipage} & \begin{minipage}[b]{\linewidth}\raggedright
Final responsibility
\end{minipage} & \begin{minipage}[b]{\linewidth}\raggedright
Implementation evidence
\end{minipage} \\
\midrule\noalign{}
\endhead
\bottomrule\noalign{}
\endlastfoot
React/Vite frontend & Candidate and HR user interface, built to static
assets & \texttt{frontend/\allowbreak{}package.\allowbreak{}json},
\texttt{scripts/\allowbreak{}ci/\allowbreak{}deploy-\allowbreak{}frontend.\allowbreak{}sh} \\
CloudFront & HTTPS public entry point and route dispatcher &
\texttt{scripts/\allowbreak{}aws/\allowbreak{}ensure-\allowbreak{}cloudfront.\allowbreak{}sh}, Week 8 distribution
\texttt{EQIGYNECXDYL8} \\
Frontend S3 bucket & Private storage for \texttt{frontend/\allowbreak{}dist} assets &
Private frontend bucket from the Week 8 deployment evidence; account
suffix redacted \\
Application Load Balancer & Public routing to EKS services &
\texttt{k8s/\allowbreak{}eks/\allowbreak{}ingress-\allowbreak{}alb-\allowbreak{}no-\allowbreak{}domain.\allowbreak{}yaml} \\
FastAPI backend & Auth, jobs, applications, uploads, dashboards,
processing job APIs & \texttt{backend/\allowbreak{}app/\allowbreak{}main.\allowbreak{}py}, backend routers \\
Chat service & REST chat APIs and Socket.IO realtime channel &
\texttt{chat-\allowbreak{}service/\allowbreak{}server.\allowbreak{}js} \\
Processing worker & Claims asynchronous processing jobs and writes
results & \texttt{backend/\allowbreak{}app/\allowbreak{}workers/\allowbreak{}processing\_\allowbreak{}worker.\allowbreak{}py},
\texttt{k8s/\allowbreak{}app/\allowbreak{}backend-\allowbreak{}processing-\allowbreak{}worker.\allowbreak{}yaml} \\
Outbox dispatcher & Publishes committed PostgreSQL outbox events to SQS
& \texttt{backend/\allowbreak{}app/\allowbreak{}workers/\allowbreak{}outbox\_\allowbreak{}dispatcher.\allowbreak{}py}, ADR-001 \\
AI service & Stable worker-facing adapter for SageMaker &
\texttt{ai\_\allowbreak{}service/\allowbreak{}app.\allowbreak{}py}, \texttt{k8s/\allowbreak{}app/\allowbreak{}ai-\allowbreak{}service.\allowbreak{}yaml} \\
RDS PostgreSQL & Transactional business data and reliable worker queues
& Alembic migrations through
\texttt{0008\_\allowbreak{}async\_\allowbreak{}processing\_\allowbreak{}jobs.\allowbreak{}py} \\
DynamoDB & Chat persistence and Lambda event deduplication & Chat table
names in Kubernetes config and Week 8 runtime evidence \\
Redis & Socket.IO pub/sub between chat pods &
\texttt{chat-\allowbreak{}service/\allowbreak{}lib/\allowbreak{}redis.\allowbreak{}js}, Kubernetes config \\
SQS and DLQ & At-least-once event delivery and failed-message isolation
& ADR-001 and Week 8 queue evidence \\
Lambda and SES & Event notification, archive, and email delivery &
Lambda configuration and CloudWatch evidence from Week 8 \\
GitHub Actions OIDC & CI/CD without long-lived AWS access keys &
\texttt{.\allowbreak{}github/\allowbreak{}workflows/\allowbreak{}cicd.\allowbreak{}yml} \\
\end{longtable}
\endgroup
}

\vspace{0.5em}\noindent\textbf{Design rationale}\par

The frontend was moved from EKS to S3 and CloudFront because it is
static after build time. This reduces Kubernetes workload count, removes
the need for a frontend Deployment, Service, HPA and PDB, and lets
CloudFront handle caching and SPA delivery.

The backend and chat service remain in EKS because they are long-running
APIs with health probes, replicas, HPA, PDB, and rolling deployment
needs. The processing worker also remains in EKS because CV parsing, job
parsing, and matching jobs can run longer than a short Lambda-style task
and need queue leases, retries, and controlled concurrency.

The AI service isolates SageMaker-specific inference logic from the
worker contract. The worker continues to call stable routes such as
\texttt{/\allowbreak{}parse-\allowbreak{}job}, \texttt{/\allowbreak{}parse-\allowbreak{}cv}, and
\texttt{/\allowbreak{}match-\allowbreak{}applications}, while the AI adapter handles SageMaker
endpoint invocation, timeout, retry, and response normalization.

The transactional outbox is used because committing application data and
publishing an event are separate failure domains. Events are first
inserted in PostgreSQL in the same transaction as the business mutation.
A dispatcher later sends them to SQS. This provides at-least-once
delivery, and Lambda uses DynamoDB conditional writes to deduplicate
\texttt{eventId}.

\vspace{0.5em}\noindent\textbf{Proposal vs final
implementation}\par

{
\begingroup
\small
\setlength{\tabcolsep}{3pt}
\renewcommand{\arraystretch}{1.15}
\begin{longtable}{@{}
  >{\raggedright\arraybackslash}p{(\linewidth - 4\tabcolsep) * \real{0.3333}}
  >{\raggedright\arraybackslash}p{(\linewidth - 4\tabcolsep) * \real{0.3333}}
  >{\raggedright\arraybackslash}p{(\linewidth - 4\tabcolsep) * \real{0.3333}}@{}}
\toprule\noalign{}
\begin{minipage}[b]{\linewidth}\raggedright
Area
\end{minipage} & \begin{minipage}[b]{\linewidth}\raggedright
Original proposal
\end{minipage} & \begin{minipage}[b]{\linewidth}\raggedright
Final implemented architecture
\end{minipage} \\
\midrule\noalign{}
\endhead
\bottomrule\noalign{}
\endlastfoot
Frontend hosting & Could be served through Kubernetes or AWS static
hosting & Implemented as private S3 plus CloudFront \\
Public entry point & ALB and/or CloudFront expected & CloudFront is the
user-facing entry point; ALB is an origin for API/chat/socket paths \\
Backend runtime & Kubernetes on EKS & EKS Deployment \texttt{backend}, 2
replicas, HPA 2-5 \\
Chat runtime & Node.js and Socket.IO with Redis/DynamoDB & EKS
Deployment \texttt{chat-\allowbreak{}service}, 2 replicas, ALB stickiness, Redis
pub/sub \\
AI runtime & AI service expected & EKS \texttt{ai-\allowbreak{}service} adapter
invokes SageMaker endpoint \texttt{internship-\allowbreak{}qwen3-\allowbreak{}4b} \\
Event processing & Asynchronous events expected & PostgreSQL outbox,
SQS, Lambda, DynamoDB dedupe, S3 archive and SES \\
Deployment & CI/CD expected & GitHub Actions workflow dispatch modes
with OIDC and ECR image verification \\
Runtime proof & Design-stage assumption & I collected runtime evidence
for CloudFront, EKS workloads, RDS, Redis, DynamoDB, SQS, Lambda
configuration, and the SageMaker endpoint \\
\end{longtable}
\endgroup
}

\vspace{0.5em}\noindent\textbf{AWS Services Used}\par

{
\begingroup
\small
\setlength{\tabcolsep}{3pt}
\renewcommand{\arraystretch}{1.15}
\begin{longtable}{@{}
  >{\raggedright\arraybackslash}p{(\linewidth - 4\tabcolsep) * \real{0.3333}}
  >{\raggedright\arraybackslash}p{(\linewidth - 4\tabcolsep) * \real{0.3333}}
  >{\raggedright\arraybackslash}p{(\linewidth - 4\tabcolsep) * \real{0.3333}}@{}}
\toprule\noalign{}
\begin{minipage}[b]{\linewidth}\raggedright
AWS service
\end{minipage} & \begin{minipage}[b]{\linewidth}\raggedright
Role in the system
\end{minipage} & \begin{minipage}[b]{\linewidth}\raggedright
Reason for selection
\end{minipage} \\
\midrule\noalign{}
\endhead
\bottomrule\noalign{}
\endlastfoot
Amazon EKS & Runs the backend, worker, chat service, and container
workloads & Supports Kubernetes, scaling, rolling updates, and
multi-service management \\
Amazon ECR & Stores container images & Integrates directly with EKS and
GitHub Actions \\
Application Load Balancer & Receives and routes requests to services &
Supports routing, health checks, and Kubernetes Ingress integration \\
Amazon RDS for PostgreSQL & Stores users, companies, jobs, applications,
document metadata, and processing jobs & Suitable for relational data
and transactions \\
Amazon DynamoDB & Stores chat users, conversations, and messages &
Suitable for chat data with high read/write traffic \\
Amazon ElastiCache & Provides Redis pub/sub for Socket.IO and temporary
data & Helps multiple chat-service instances synchronize events \\
Amazon S3 & Stores CVs, certificates, transcripts, and frontend assets &
Provides scalability, high durability, and presigned URL support \\
Amazon CloudFront & Distributes frontend and static assets & Improves
access speed and supports HTTPS \\
Amazon SQS & Receives asynchronous events from the outbox dispatcher &
Reduces direct coupling between services \\
Amazon SageMaker & Provides an AI inference endpoint in production &
Supports managed model endpoints and scalability \\
Amazon CloudWatch & Collects AWS logs, metrics, and alarms & Supports
operations and incident investigation \\
AWS IAM & Manages access between workloads and AWS resources & Supports
the least privilege principle \\
\end{longtable}
\endgroup
}

\vspace{0.5em}\noindent\textbf{Component Design}\par

\textbf{Frontend}

The frontend is developed with React and Vite and provides interfaces
for Candidates and HR. Static files can be built and stored on Amazon
S3, then distributed through Amazon CloudFront. The frontend
communicates with the backend and chat service through endpoints exposed
by the Application Load Balancer.

\textbf{Backend API}

The backend is developed with FastAPI and is responsible for user
authentication, Candidate and HR management, company and job management,
application and document metadata management, presigned URL generation,
processing job creation, dashboard data, and business APIs.

\textbf{Processing Worker}

The processing worker performs long-running tasks such as CV content
extraction, Job Description analysis, CV analysis, matching score
calculation, and candidate reranking. Lease, retry, and attempt-limit
mechanisms are used to reduce failures when a worker stops during
processing.

\textbf{Chat Service}

The chat service uses Node.js, Express, and Socket.IO to support
realtime communication. DynamoDB stores users, groups, and messages; the
Redis adapter synchronizes events across multiple pods; sticky sessions
help maintain Socket.IO connections when multiple replicas are running.

\textbf{AI Service}

The AI service analyzes CVs and Job Descriptions, converts text into
structured data, and supports fit scoring. The AI service should include
schema validation, timeout, limited retry, health checks, log control to
avoid exposing CV data, and fallback behavior when the model is
unavailable.

\textbf{Data Storage}

PostgreSQL stores relational data that requires transactions. DynamoDB
stores chat data. Amazon S3 stores large files. Redis supports realtime
event distribution. Using multiple storage technologies allows each data
group to be managed with a technology that matches its access pattern.

\textbf{Observability}

The system collects metrics for performance and health monitoring, logs
for event and error investigation, and distributed traces to follow
requests across services. Prometheus, Grafana, Loki, OpenTelemetry, and
Tempo are used in Kubernetes environments; Amazon CloudWatch is used for
AWS resource and log monitoring.

\vspace{0.5em}\noindent\textbf{Technical Implementation}\par

\vspace{0.5em}\noindent\textbf{Implementation Phases}\par

{
\begingroup
\small
\setlength{\tabcolsep}{3pt}
\renewcommand{\arraystretch}{1.15}
\begin{longtable}{@{}
  >{\raggedright\arraybackslash}p{(\linewidth - 2\tabcolsep) * \real{0.5000}}
  >{\raggedright\arraybackslash}p{(\linewidth - 2\tabcolsep) * \real{0.5000}}@{}}
\toprule\noalign{}
\begin{minipage}[b]{\linewidth}\raggedright
Phase
\end{minipage} & \begin{minipage}[b]{\linewidth}\raggedright
Main activities
\end{minipage} \\
\midrule\noalign{}
\endhead
\bottomrule\noalign{}
\endlastfoot
1. Analysis and design & Analyze the Candidate-HR problem, define
functional and non-functional requirements, design the database, design
the multi-service architecture, and identify required AWS services \\
2. Business platform foundation & Build authentication, Candidate/HR
profiles, company management, job management, application tracking, and
document management \\
3. Realtime and AI integration & Build the chat service, integrate
Socket.IO, store chat data in DynamoDB, integrate Redis adapter, build
the AI service, analyze CV/JD content, and implement matching and
reranking \\
4. Reliability improvements & Build processing workers, retry, lease,
idempotency, optimistic concurrency, transactional outbox, validation,
and error handling \\
5. Containerization and Kubernetes & Write Dockerfiles, build Docker
Compose, create a local Kubernetes cluster with kind, deploy Deployment,
Service, Ingress, HPA, PDB, health probes, and migration jobs \\
6. Observability and CI/CD & Collect metrics, logs, and traces, build
dashboards, set up alerts, build GitHub Actions, run automated tests,
smoke tests, and security scans \\
7. AWS deployment and completion & Build images, push to Amazon ECR,
deploy workloads on Amazon EKS, connect RDS, DynamoDB, ElastiCache, and
S3, deploy frontend on S3/CloudFront, run end-to-end tests, and complete
the documentation \\
\end{longtable}
\endgroup
}

\vspace{0.5em}\noindent\textbf{Technical Requirements}\par

{
\begingroup
\small
\setlength{\tabcolsep}{3pt}
\renewcommand{\arraystretch}{1.15}
\begin{longtable}{@{}
  >{\raggedright\arraybackslash}p{(\linewidth - 2\tabcolsep) * \real{0.5000}}
  >{\raggedright\arraybackslash}p{(\linewidth - 2\tabcolsep) * \real{0.5000}}@{}}
\toprule\noalign{}
\begin{minipage}[b]{\linewidth}\raggedright
Component
\end{minipage} & \begin{minipage}[b]{\linewidth}\raggedright
Technology or requirement
\end{minipage} \\
\midrule\noalign{}
\endhead
\bottomrule\noalign{}
\endlastfoot
Frontend & React, Vite, Tailwind CSS \\
Backend & Python, FastAPI, SQLAlchemy, Alembic \\
Chat service & Node.js, Express, Socket.IO \\
AI service & Python, FastAPI, Qwen-compatible model or SageMaker
endpoint \\
Relational database & PostgreSQL \\
Chat database & DynamoDB \\
Cache/pub-sub & Redis or Amazon ElastiCache \\
Object storage & Amazon S3 \\
Containers & Docker \\
Container orchestration & Kubernetes, kind, and Amazon EKS \\
CI/CD & GitHub Actions \\
Monitoring & Prometheus, Grafana, and CloudWatch \\
Logs & Loki and CloudWatch Logs \\
Tracing & OpenTelemetry and Tempo \\
Security & JWT, IAM Role, OIDC, least privilege, and secret
management \\
\end{longtable}
\endgroup
}

\vspace{0.5em}\noindent\textbf{Timeline and Milestones}\par

I adjusted the plan to 8 weeks, from 08/06/2026 to 30/07/2026, to match
the worklog timeline.

{
\begingroup
\small
\setlength{\tabcolsep}{3pt}
\renewcommand{\arraystretch}{1.15}
\begin{longtable}{@{}
  >{\raggedright\arraybackslash}p{(\linewidth - 6\tabcolsep) * \real{0.2500}}
  >{\raggedright\arraybackslash}p{(\linewidth - 6\tabcolsep) * \real{0.2500}}
  >{\raggedright\arraybackslash}p{(\linewidth - 6\tabcolsep) * \real{0.2500}}
  >{\raggedright\arraybackslash}p{(\linewidth - 6\tabcolsep) * \real{0.2500}}@{}}
\toprule\noalign{}
\begin{minipage}[b]{\linewidth}\raggedright
Week
\end{minipage} & \begin{minipage}[b]{\linewidth}\raggedright
Time
\end{minipage} & \begin{minipage}[b]{\linewidth}\raggedright
Milestone
\end{minipage} & \begin{minipage}[b]{\linewidth}\raggedright
Expected deliverable
\end{minipage} \\
\midrule\noalign{}
\endhead
\bottomrule\noalign{}
\endlastfoot
1 & 08/06/2026 - 14/06/2026 & Requirements analysis, architecture,
backend foundation, and authentication & Project scope, high-level
architecture, registration, login, and initial migration \\
2 & 15/06/2026 - 21/06/2026 & Candidate-HR platform, applications,
documents, and S3 & Company, jobs, application flow, secure upload, and
document download \\
3 & 22/06/2026 - 28/06/2026 & Frontend and REST API integration &
Candidate/HR interfaces, protected routes, and API integration \\
4 & 29/06/2026 - 05/07/2026 & Realtime chat & Realtime communication
between Candidate and HR using Socket.IO, Redis, and DynamoDB \\
5 & 06/07/2026 - 12/07/2026 & AI integration and reliable processing &
CV/JD analysis, matching, worker, idempotency, concurrency, and
outbox \\
6 & 13/07/2026 - 19/07/2026 & Docker Compose & Full stack runs locally
with smoke test procedure \\
7 & 20/07/2026 - 26/07/2026 & Kubernetes and observability & System runs
on kind with metrics, logs, traces, and alerts \\
8 & 27/07/2026 - 30/07/2026 & CI/CD, AWS, and reporting & End-to-end
testing, AWS deployment preparation, and completed documentation \\
\end{longtable}
\endgroup
}

\vspace{0.5em}\noindent\textbf{Budget Estimation}\par

Exact monthly cost should be calculated with AWS Pricing Calculator and
then compared with representative billing data for the project AWS
account in \texttt{ap-\allowbreak{}southeast-\allowbreak{}1}. For this report, the Week 8 AWS
evidence is treated as the source of truth for current cost status.

\vspace{0.5em}\noindent\textbf{Week 8 AWS cost evidence}\par

The Week 8 evidence reports July 1-28, 2026 month-to-date AWS spend and
a Billing and Cost Management credits screenshot. The finalized monthly
bill and AWS Pricing Calculator estimate are still pending.

{
\begingroup
\small
\setlength{\tabcolsep}{3pt}
\renewcommand{\arraystretch}{1.15}
\begin{longtable}{@{}
  >{\raggedright\arraybackslash}p{(\linewidth - 4\tabcolsep) * \real{0.3000}}
  >{\raggedleft\arraybackslash}p{(\linewidth - 4\tabcolsep) * \real{0.4000}}
  >{\raggedright\arraybackslash}p{(\linewidth - 4\tabcolsep) * \real{0.3000}}@{}}
\toprule\noalign{}
\begin{minipage}[b]{\linewidth}\raggedright
Cost evidence
\end{minipage} & \begin{minipage}[b]{\linewidth}\raggedleft
Observed value
\end{minipage} & \begin{minipage}[b]{\linewidth}\raggedright
Interpretation
\end{minipage} \\
\midrule\noalign{}
\endhead
\bottomrule\noalign{}
\endlastfoot
July 1-28, 2026 total spend & \texttt{\$94.\allowbreak{}92} & Current month-to-date
spend from the Week 8 cost summary \\
Highest daily spend & \texttt{\$31.\allowbreak{}83} on July 28 & Largest daily cost
in the July 1-28 period \\
Credits total amount used & \texttt{\$27.\allowbreak{}90} & Billing credits
screenshot evidence \\
Credits total estimated amount used & \texttt{\$140.\allowbreak{}65} & Billing
credits screenshot evidence \\
Credits total amount remaining & \texttt{\$172.\allowbreak{}10} & Billing credits
screenshot evidence \\
Credits total estimated amount remaining & \texttt{\$59.\allowbreak{}35} & Billing
credits screenshot evidence \\
\end{longtable}
\endgroup
}

The top Week 8 cost drivers are Amazon RDS (\texttt{\$29.\allowbreak{}69}), Amazon
SageMaker (\texttt{\$23.\allowbreak{}45}), Amazon VPC (\texttt{\$14.\allowbreak{}11}), Amazon EC2
- Compute (\texttt{\$12.\allowbreak{}42}), EC2 - Other (\texttt{\$7.\allowbreak{}68}), and Amazon
EKS (\texttt{\$5.\allowbreak{}64}). Together they account for approximately 98\% of
the reported July 1-28 spend.

\vspace{0.5em}\noindent\textbf{Cost assumptions}\par

{
\begingroup
\small
\setlength{\tabcolsep}{3pt}
\renewcommand{\arraystretch}{1.15}
\begin{longtable}{@{}
  >{\raggedright\arraybackslash}p{(\linewidth - 4\tabcolsep) * \real{0.3333}}
  >{\raggedright\arraybackslash}p{(\linewidth - 4\tabcolsep) * \real{0.3333}}
  >{\raggedright\arraybackslash}p{(\linewidth - 4\tabcolsep) * \real{0.3333}}@{}}
\toprule\noalign{}
\begin{minipage}[b]{\linewidth}\raggedright
Area
\end{minipage} & \begin{minipage}[b]{\linewidth}\raggedright
Evidence-based assumption
\end{minipage} & \begin{minipage}[b]{\linewidth}\raggedright
Estimate status
\end{minipage} \\
\midrule\noalign{}
\endhead
\bottomrule\noalign{}
\endlastfoot
Environment & Production deployment in \texttt{ap-\allowbreak{}southeast-\allowbreak{}1} &
Verified \\
Public traffic & CloudFront fronts browser traffic and routes API/chat
traffic to ALB & Verified \\
Frontend & Static assets are stored in a private S3 bucket and delivered
by CloudFront & Verified \\
Kubernetes & One EKS cluster runs backend, chat, worker, dispatcher, and
AI adapter workloads & Verified \\
Worker nodes & Two Ready Kubernetes worker nodes were captured in the
evidence snapshot & Instance type and EBS size still required \\
RDS PostgreSQL & Two private, encrypted PostgreSQL \texttt{db.\allowbreak{}t4g.\allowbreak{}micro}
instances with 20 GiB storage each & Verified \\
ElastiCache / Valkey & One available Redis-compatible replication group
with encryption at rest and in transit & Node type/count still
required \\
DynamoDB & Chat and Lambda dedupe tables use on-demand billing and are
near empty in the snapshot & Verified \\
SQS & Main outbox queue and DLQ have zero visible messages in the
snapshot & Verified \\
Lambda & Outbox handler is Active, 256 MB memory, 20 second timeout &
Verified \\
SageMaker & Real-time endpoint is \texttt{InService} with one production
variant & Instance type and expected uptime still required \\
CloudWatch & Logs and metrics exist through AWS and Kubernetes services
& Retention and ingestion volume still required \\
\end{longtable}
\endgroup
}

\vspace{0.5em}\noindent\textbf{Cost-estimation
methodology}\par

\begin{enumerate}
\def\labelenumi{\arabic{enumi}.}
\tightlist
\item
  Use AWS Pricing Calculator for \texttt{ap-\allowbreak{}southeast-\allowbreak{}1}.
\item
  Enter the observed deployed resources from the evidence folder.
\item
  Add missing runtime inputs: node instance type, EBS size, Redis node
  type/count, SageMaker instance type, monthly runtime hours, data
  transfer, request volume, and CloudWatch log retention.
\item
  Export the Pricing Calculator estimate as the forecast.
\item
  Compare the forecast with Cost Explorer after the environment has run
  for a representative period.
\end{enumerate}

{
\begingroup
\small
\setlength{\tabcolsep}{3pt}
\renewcommand{\arraystretch}{1.15}
\begin{longtable}{@{}
  >{\raggedright\arraybackslash}p{(\linewidth - 4\tabcolsep) * \real{0.3333}}
  >{\raggedright\arraybackslash}p{(\linewidth - 4\tabcolsep) * \real{0.3333}}
  >{\raggedright\arraybackslash}p{(\linewidth - 4\tabcolsep) * \real{0.3333}}@{}}
\toprule\noalign{}
\begin{minipage}[b]{\linewidth}\raggedright
Service
\end{minipage} & \begin{minipage}[b]{\linewidth}\raggedright
Monthly estimate method
\end{minipage} & \begin{minipage}[b]{\linewidth}\raggedright
Current evidence status
\end{minipage} \\
\midrule\noalign{}
\endhead
\bottomrule\noalign{}
\endlastfoot
Amazon EKS & Cluster hourly price times monthly runtime & Deployed; Week
8 cost summary reports \texttt{\$5.\allowbreak{}64} \\
EC2 worker nodes & Node hourly price times two nodes, plus EBS storage &
Nodes verified; EC2 Compute reports \texttt{\$12.\allowbreak{}42}, EC2 - Other
reports \texttt{\$7.\allowbreak{}68} \\
Amazon RDS PostgreSQL & \texttt{db.\allowbreak{}t4g.\allowbreak{}micro} hourly price, 20 GiB
storage per instance, backup/I/O & Two private encrypted instances
verified; Week 8 cost summary reports \texttt{\$29.\allowbreak{}69} \\
Amazon ElastiCache / Valkey & Node hourly price times node count, plus
data transfer & Replication group verified; node type/count needed \\
Application Load Balancer & ALB hourly charge plus LCU usage & Active
ALB verified; included in networking/runtime cost checks \\
Amazon CloudFront & Requests, transfer out, and invalidations &
Monitoring screenshot available; behavior/origin export still
required \\
Amazon S3 & Storage GB, PUT/GET requests, lifecycle transitions &
Frontend bucket object screenshot available \\
Amazon DynamoDB & On-demand read/write requests plus storage & Table
status screenshot available \\
Amazon SQS & Standard queue requests and payload volume & Queue and DLQ
screenshot available \\
AWS Lambda & Requests plus GB-seconds & Function exists; invocation cost
and trigger health still pending \\
Amazon SES & Email send count and attachments if any & Used by
notification path \\
Amazon ECR & Image storage and transfer if applicable & Not separated in
the Week 8 top-driver table \\
Amazon SageMaker & Endpoint instance hours and invocation/data charges &
Endpoint verified \texttt{InService}; Week 8 cost summary reports
\texttt{\$23.\allowbreak{}45} \\
Amazon CloudWatch & Log ingestion, storage retention, metrics, alarms &
Retention and ingestion volume still required \\
Amazon VPC and data transfer & CloudFront egress, ALB traffic, NAT
processing, cross-AZ traffic & Week 8 cost summary reports Amazon VPC at
\texttt{\$14.\allowbreak{}11}; traffic breakdown still needed \\
Total & Pricing Calculator forecast, then billing comparison & Week 8
July 1-28 spend reports \texttt{\$94.\allowbreak{}92}; finalized bill and
steady-state estimate pending \\
\end{longtable}
\endgroup
}

The largest current and future cost drivers are RDS, SageMaker endpoint
uptime, VPC/NAT-related networking, EC2 worker nodes and EBS, EKS
control plane, Redis, ALB, CloudFront/data transfer, and CloudWatch
logs. If the SageMaker endpoint remains online continuously, AI
inference can become one of the dominant monthly costs.

\vspace{0.5em}\noindent\textbf{Evidence hygiene for cost
materials}\par

No raw screenshots or logs should expose account IDs, access keys,
secret keys, database passwords, connection strings, Kubernetes Secret
values, GitHub tokens, presigned URLs, or login-related logs. The pod
log evidence must be summarized or sanitized before publication because
the raw \texttt{02-\allowbreak{}eks/\allowbreak{}pod-\allowbreak{}logs-\allowbreak{}tail.\allowbreak{}txt} file contains a Redis
connection URL with an embedded credential.

\vspace{0.5em}\noindent\textbf{Cost Optimization}\par

\begin{itemize}
\tightlist
\item
  Run the production environment only when demoing or testing.
\item
  Use instance types that match the actual workload.
\item
  Limit the number of worker nodes.
\item
  Stop or delete SageMaker endpoints when not in use.
\item
  Configure suitable log retention.
\item
  Delete old images in ECR.
\item
  Apply S3 lifecycle policies when appropriate.
\item
  Use DynamoDB on-demand while traffic is still unstable.
\item
  Configure AWS Budget and Billing Alarm.
\item
  Delete unused Load Balancers, public IPv4 addresses, snapshots, and
  volumes.
\item
  Confirm whether NAT Gateway is required before leaving it running.
\end{itemize}

\vspace{0.5em}\noindent\textbf{Risk Assessment}\par

\vspace{0.5em}\noindent\textbf{Risk Matrix}\par

{
\begingroup
\small
\setlength{\tabcolsep}{3pt}
\renewcommand{\arraystretch}{1.15}
\begin{longtable}{@{}
  >{\raggedright\arraybackslash}p{(\linewidth - 6\tabcolsep) * \real{0.2500}}
  >{\raggedright\arraybackslash}p{(\linewidth - 6\tabcolsep) * \real{0.2500}}
  >{\raggedright\arraybackslash}p{(\linewidth - 6\tabcolsep) * \real{0.2500}}
  >{\raggedright\arraybackslash}p{(\linewidth - 6\tabcolsep) * \real{0.2500}}@{}}
\toprule\noalign{}
\begin{minipage}[b]{\linewidth}\raggedright
Risk
\end{minipage} & \begin{minipage}[b]{\linewidth}\raggedright
Likelihood
\end{minipage} & \begin{minipage}[b]{\linewidth}\raggedright
Impact
\end{minipage} & \begin{minipage}[b]{\linewidth}\raggedright
Severity
\end{minipage} \\
\midrule\noalign{}
\endhead
\bottomrule\noalign{}
\endlastfoot
Secret or AWS credential exposure & Low & Very high & High \\
RDS or services become public unnecessarily & Medium & High & High \\
CV and personal data are accessed without permission & Low & Very high &
High \\
AI returns inaccurate results & Medium & Medium & Medium \\
AI service runs out of resources or times out & Medium & High & High \\
Realtime chat disconnects when scaled & Medium & Medium & Medium \\
Worker processes duplicate jobs & Medium & High & High \\
Database updates encounter race conditions & Medium & High & High \\
Kubernetes pod or node failure & Medium & Medium & Medium \\
AWS cost exceeds expectations & Medium & High & High \\
CI/CD deployment fails & Medium & Medium & Medium \\
Third-party package has a vulnerability & Medium & High & High \\
\end{longtable}
\endgroup
}

\vspace{0.5em}\noindent\textbf{Mitigation Measures}\par

\textbf{Security}

\begin{itemize}
\tightlist
\item
  Use IAM Roles and least privilege.
\item
  Do not store AWS Access Keys in source code.
\item
  Use GitHub Actions OIDC.
\item
  Keep \texttt{.\allowbreak{}env}, private keys, and secrets out of the repository.
\item
  Keep S3 buckets private.
\item
  Use presigned URLs with expiration.
\item
  Do not log full CV content or tokens.
\item
  Check authorization and object ownership in the backend.
\end{itemize}

\textbf{Reliability}

\begin{itemize}
\tightlist
\item
  Use readiness and liveness probes.
\item
  Use HPA to support scaling.
\item
  Use PDB to reduce disruption during maintenance.
\item
  Use limited retry.
\item
  Apply processing-job leases.
\item
  Use idempotency keys.
\item
  Apply optimistic concurrency.
\item
  Use transactional outbox for events.
\end{itemize}

\textbf{AI}

\begin{itemize}
\tightlist
\item
  Validate output with schemas.
\item
  Set timeouts.
\item
  Limit retries.
\item
  Allow human review before recruitment decisions.
\item
  Do not use AI scores as the only hiring decision.
\item
  Record reasons or contributing factors when possible.
\end{itemize}

\textbf{Cost}

\begin{itemize}
\tightlist
\item
  Configure AWS Budget.
\item
  Monitor Cost Explorer.
\item
  Delete resources after the workshop.
\item
  Limit AI endpoint runtime.
\item
  Use suitable log retention.
\item
  Check for remaining EBS volumes, Elastic IPs, Load Balancers, and
  snapshots.
\end{itemize}

\vspace{0.5em}\noindent\textbf{Contingency Plan}\par

\begin{itemize}
\tightlist
\item
  If EKS is not ready, use Docker Compose or local Kubernetes for the
  demo.
\item
  If the AI model is unavailable, use mock responses or deterministic
  scoring to continue testing the business flow.
\item
  If Redis is unavailable, run a single chat-service instance for the
  demo environment.
\item
  If S3 cannot be accessed, use local storage in development.
\item
  If a worker fails, processing jobs are retried with lease and attempt
  limits.
\item
  If a new deployment fails, roll back to the previous container image
  or Kubernetes revision.
\item
  If cost exceeds the limit, stop non-essential workloads and delete
  high-cost resources.
\end{itemize}

\vspace{0.5em}\noindent\textbf{Expected Outcomes}\par

After completion, the system is expected to achieve the following
outcomes:

\begin{itemize}
\tightlist
\item
  Provide a centralized platform for Candidates and HR.
\item
  Manage companies, jobs, applications, and documents.
\item
  Store CVs and documents securely.
\item
  Support realtime chat.
\item
  Support AI analysis for CVs and Job Descriptions.
\item
  Process long-running tasks with workers.
\item
  Control repeated requests and concurrent updates.
\item
  Run with Docker Compose and Kubernetes.
\item
  Provide a CI/CD pipeline.
\item
  Provide metrics, logs, traces, and alerts.
\item
  Provide an AWS deployment architecture.
\item
  Provide documentation for deployment, testing, and cleanup.
\end{itemize}

\vspace{0.5em}\noindent\textbf{Long-Term Value}\par

The project can continue evolving into a more complete system with:

\begin{itemize}
\tightlist
\item
  Email or notification when application status changes.
\item
  Interview scheduling.
\item
  AI-generated cover letters.
\item
  Interview questions generated from CVs and Job Descriptions.
\item
  Job recommendations based on Candidate profiles.
\item
  Recruitment analytics dashboard.
\item
  Calendar and email integration.
\item
  Improved explainability for AI matching.
\item
  Mobile application support.
\item
  Expansion to multiple universities and companies.
\end{itemize}

Beyond product value, the project also provides a reference for building
a Cloud-native multi-service system with AI integration, realtime
communication, asynchronous processing, Kubernetes, observability, and
CI/CD.
